% ===== §14 The Ethical Limits of Aesthetic Education =====
\section{The Ethical Limits of Aesthetic Education}
\label{sec:limits}

\noindent
A paper that has proposed an education must say where the education stops, and it must do so for reasons internal to what the education is rather than as a concession appended at the end. The cultivation of perception has limits that are not external boundaries fencing in an otherwise unlimited practice but conditions of its being the practice it is; to cross them is not to do too much of a good thing but to convert the practice into its opposite. This section states three such limits. The first is that the generativity aesthetic education cultivates cannot be made procedural without being killed. The second is that the cultivation of perception must be told apart from its disciplinary double, the training of a docile perception, and that the two share a mechanism and differ in what they produce and to whom its value returns. The third the section states not as a proposition but by instantiating it, in a return to the ordinary life from which the paper set out.

\subsection{Against procedure}
\label{sub:limits-procedure}

The first limit was implicit in the paper's method from the start and can now be made explicit. The generativity that aesthetic education cultivates is killed by being made procedural. A generation that could be produced by following a fixed procedure would not be a generation but an execution, the enactment of what was laid down in advance, and the field of contingency in which alone perception generates would have been replaced by an arrangement in which it is merely performed (\xref{sub:diag-contingency}). This is why the paper has declined throughout to convert its models into techniques, and why the fourth model made the openness of aesthetic practice a structural requirement rather than a modesty: if the good aesthetic arrangement is not presupposed but emerges through open-ended refinement, then any procedure claiming to produce it in advance has mistaken the closed cycle for the open one (\xref{sub:typ-model4}). The models offered are therefore not procedures and cannot be made into them without ceasing to be what they are. They are descriptions of the structure of a generative practice, offered to be inhabited and refined, not steps to be followed to a guaranteed result. An aesthetic education administered as a procedure, a curriculum of perception with prescribed inputs and examinable outputs, would produce exactly the settled and appropriable perception the paper has opposed at every register, and would do so the more insidiously for calling itself an education of generativity. The limit is strict: to the precise extent that the cultivation of perception is proceduralised, it stops cultivating generativity and begins producing its counterfeit.

\subsection{Cultivation and discipline}
\label{sub:limits-discipline}

The first limit opens directly onto the second, because the proceduralised training of perception has a name and a theorist, and confronting it is how the paper draws the line between what it proposes and what it must refuse. Foucault described discipline as the production of docile bodies, a training that works on the body through exercise and repetition to produce a subject that is useful, obedient, and improved in the terms of the power that trains it, and that internalises the norm to which it is trained until surveillance need no longer be external \citep{foucault1977}. Discipline is a cultivation of the body, and it works, as the cultivation of perception works, through the plasticity that makes the body trainable and through the repeated practice that reshapes it. The two cannot be told apart by their mechanism, for the mechanism is one: both retune a plastic perceptual and bodily capacity through sustained training. They are told apart by what they produce and to whom its value returns, which is the criterion the paper has carried throughout. Discipline produces docility: a perception trained to the norm, tuned to obey and to be useful to the power that trained it, its value extracted to that power and the subject preserved only as an instrument, the negative holonomy of a training whose surplus is drawn off from the one trained. Cultivation, in the sense this paper has given it, produces generativity: a perception trained toward its own enlargement, whose value returns to the perceiver and to the shared attending that generated it, and in which the subject is preserved as a generator rather than reduced to an instrument, the positive holonomy of a training whose surplus returns to the one trained. The distinction is exactly the one the theory of relational reproduction fixed in its principle of the preservation of the subject (\xref{sub:typ-model1}), and the one the second spine drew between refinement and commodified training (\xref{sub:plas-redline}). Aesthetic education stands on one side of this line, and the disciplinary training of a docile perception on the other, and because they share a mechanism the line must be drawn with care and held with vigilance, for an aesthetic education that forgot the distinction would slide, under the pressure to produce results, into the disciplinary production of a perception trained to consume, obey, and be useful, which is precisely the perception the spectacle already produces at scale (\xref{sub:dist-spectacle}). The limit is therefore not only a caution but a criterion of self-examination: an aesthetic education must ask, of its every practice, whether it enlarges the generativity of those it cultivates and returns to them the value of their trained perception, or whether it has begun, however well-meaningly, to produce a docility whose value it extracts.

\subsection{The ordinary evening}
\label{sub:limits-instantiation}

The third limit the paper states by instantiating it, since a paper on the generative perception of ordinary life should be able to return, at its close, to an ordinary evening and find there what it has been theorising. Consider, then, without further abstraction, two people at the end of a day. Neither is the teacher of the other, and neither is the other's pupil; what passes between them is not a lesson in perception but the mutual and unforced training of it, each drawing the other's attention, without design, to a quality of the light, a turn of phrase overheard, the particular way the evening has arranged itself, and each finding their own perception altered by what the other has taught them to see. They arrange their shared space not to a plan but by the accumulated small acts of a life made together, leaving room where room is wanted and marking with a word or a kept object what has come to matter between them. They fall silent, sometimes, and let an experience settle before returning it, judged in common rather than in the solitude of either, to the shared life it belongs to. And over the long iteration of their days they refine, without ever naming it as refinement, the strategies by which they show each other tenderness, keeping what the other's response returns as generative and letting go what does not, the good of it emerging between them rather than known in advance by either. This is the whole of the paper's proposal, instantiated in an evening that asks for no theory to occur and that occurs, unnumbered, in ordinary lives. The limit it instantiates is the one the paper most needs its reader to feel rather than to be told: that the cultivation of perception this paper has theorised is not a programme to be administered but a description of what two people, turned together toward their shared and ordinary days, already do when their common life goes well, and that the office of the theory is not to produce this but to see it, to give it a name, and to defend it against the forces that would settle, extract, and proceduralise it out of existence. The theory ends where the evening begins, and hands what it has understood back to the ordinary life it was always about.
